\section[Notebook 1982-7: symmetric matrices with non-positive off-diagonal entries]{Notebook 1982-7: a symmetric matrix with non-positive off-diagonal entries and a positive vector $v$ with $Av \ge 0$ is positive semidefinite\protect\verifie{Notebook1982_7}}\label{sec:quillen-1982-7}

\emph{Context.} \carnet{1982}{7} (11 June -- 12 July 1982, his pages 700--787) works on determinants of $\bar\partial$-operators when the complex structure of the Riemann surface varies, on Riemann--Roch and the Hodge index theorem for surfaces, and on Arakelov's arithmetic surfaces (with notes of a talk by Faltings on 30 June). In the entry of 25 June (pages~32--34) Quillen wants to understand the Hodge index theorem through Mumford's result that the intersection matrix $(C_i \cdot C_j)$ of the exceptional curves of a resolved surface singularity is negative definite. He reduces it to an algebraic situation: a symmetric matrix with entries $\ge 0$ off the diagonal and a combination $D = \sum n_i C_i$ with all $n_i > 0$ and $C_i \cdot D = 0$ for every $i$.

\emph{The page.} On page~33 he passes to $A = -(C_i \cdot C_j)$ and tries an induction on the size, splitting a vector into its positive and negative parts; he notes that subtracting a multiple of $D$ destroys the induction. On page~34 he states the correct assertion: a real symmetric matrix with off-diagonal entries $\le 0$, for which some vector $v$ with strictly positive entries satisfies the hypothesis, is positive semidefinite; after normalising $v = (1, \ldots, 1)$ the proof uses that the row sums are $\ge 0$, then lowers the diagonal entries until the row sums vanish. The hypothesis is written $v^{t}Av \ge 0$; the proof uses $Av \ge 0$, and with $v^{t}Av \ge 0$ alone the assertion is false.

\emph{Intuition.} Write $x_i = v_i y_i$. The quadratic form $x^{t}Ax$ is then $\sum_i (Av)_i v_i y_i^2$ minus half of $\sum_{i,j} A_{ij} v_i v_j (y_i - y_j)^2$: the first sum is $\ge 0$ because $Av \ge 0$, the second is $\le 0$ because only the off-diagonal entries, all $\le 0$, contribute to it.

\begin{encadre}
\begin{theoreme}[Quillen, notebook 1982-7, pp.~33--34]\label{thm:quillen-1982-7}
Let $A$ be a real symmetric $n \times n$ matrix with $A_{ij} \le 0$ for $i \ne j$.
\begin{enumerate}
\item For every vector $v$ with non-zero entries and every $x$,
\[
  x^{t}Ax = \sum_i (Av)_i\,\frac{x_i^2}{v_i} - \frac12 \sum_{i,j} A_{ij}\,v_iv_j\Bigl(\frac{x_i}{v_i} - \frac{x_j}{v_j}\Bigr)^2
\]
(this holds for any symmetric matrix over a field of characteristic $\ne 2$).
\item If some $v$ with $v_i > 0$ for all $i$ satisfies $(Av)_i \ge 0$ for all $i$, then $A$ is positive semidefinite.
\item If moreover $Av = 0$ and $x^{t}Ax = 0$, then $x_i/v_i = x_j/v_j$ whenever $A_{ij} < 0$; in particular, when the graph of the non-zero off-diagonal entries is connected, $x$ is a multiple of $v$.
\item Equivalently, for an intersection matrix $C$ with $C_{ij} \ge 0$ for $i \ne j$ and a vector $D > 0$ with $CD \le 0$, one has $x^{t}Cx \le 0$ for all $x$.
\item The hypothesis $v^{t}Av \ge 0$ does not suffice: $A = \begin{pmatrix} 4 & -1 \\ -1 & -1 \end{pmatrix}$, $v = (1, 1)$ gives $v^{t}Av = 1$, but $(0, 1)A(0, 1)^{t} = -1$.
\end{enumerate}
\end{theoreme}
\end{encadre}

\begin{proof}
For (1), put $y_i = x_i/v_i$. The left side is $\sum_{i,j} A_{ij} v_iv_j y_iy_j$, and $\sum_i (Av)_i x_i^2/v_i = \sum_{i,j} A_{ij} v_iv_j y_i^2$. By symmetry of $A$, $\sum_{i,j} A_{ij} v_iv_j y_j^2$ equals the same sum, so half of $\sum_{i,j} A_{ij} v_iv_j (y_i - y_j)^2$ is $\sum_{i,j} A_{ij}v_iv_j y_i^2 - \sum_{i,j} A_{ij} v_iv_j y_iy_j$, which is the identity. In (2) the first sum is $\ge 0$; in the double sum the diagonal terms vanish and the others are products of $A_{ij} \le 0$ with non-negative factors. In (3) the first sum is $0$, so every term of the double sum, all $\le 0$, vanishes; for $A_{ij} < 0$ this forces $y_i = y_j$. (4) is (2) applied to $-C$. (5) is a computation.
\end{proof}

\noindent The page's argument (induction on the principal submatrices, then lowering the diagonal) is not the one formalised; the identity (1) replaces it. In Quillen's application, $D = \sum n_i C_i$ is the fundamental cycle and $C_i \cdot D = 0$; part (4) gives the negative semidefiniteness he wanted, and part (3) Zariski's form of the equality case.

\emph{Verification in Lean.} The file \href{\leanblobmain/Notebook1982_7.lean}{\texttt{Notebook1982\_7.\allowbreak{}lean}} proves part~(1) as \texttt{quadratic\_identity}, over any field with $2 \ne 0$; part~(2) as \texttt{posSemidef\_of\_mulVec\_nonneg}, with mathlib's \texttt{Matrix.PosSemidef}; part~(3) as \texttt{ratio\_eq\_of\_quadratic\_eq\_zero}; part~(4) as \texttt{quadratic\_nonpos\_of\_mulVec\_eq\_zero}; and part~(5) as \texttt{not\_posSemidef\_example}. The axioms used are at most the three usual axioms of mathlib.

\litterature{the statement is classical. In the theory of $M$-matrices, a real matrix with non-positive off-diagonal entries that admits a positive vector $x$ with $Ax > 0$ is a non-singular $M$-matrix (M.~Fiedler and V.~Pták, \emph{Czechoslovak Math.\ J.} 12 (1962), 382--400); the semidefinite symmetric case is in the standard accounts of $M$-matrices. In the theory of surfaces it is Zariski's lemma on the intersection form of the components of a fibre (W.~Barth, K.~Hulek, C.~Peters and A.~Van de Ven, \emph{Compact complex surfaces}, Springer, 2nd ed.\ 2004, Lemma~III.8.2, as recalled), and the negativity of exceptional configurations is due to D.~Mumford, \emph{Publ.\ Math.\ IHÉS} 9 (1961), 5--22. The page's hypothesis $v^{t}Av \ge 0$ is a slip of the pen, as part~(5) shows; nothing here is presented as new.}
